\documentclass[11pt]{article}
\textwidth 15.9cm
\textheight 23. cm
\addtolength{\oddsidemargin}{-11.5mm}
\addtolength{\evensidemargin}{-11.5mm}
\addtolength{\topmargin}{-19.0mm}

%\usepackage{amsmath}
\usepackage[T1]{fontenc}
\usepackage[ansinew]{inputenc}
\usepackage[bbgreekl]{mathbbol}
\usepackage{amsmath,amssymb}
\usepackage{graphicx}
\usepackage[colorlinks=true,urlcolor=cyan]{hyperref}%,linkcolor=green]{hyperref}
\usepackage{url}
\usepackage[numbers,square]{natbib}
\usepackage{multirow}
\usepackage{tikz}
\usepackage{tikz-cd}
\usetikzlibrary{arrows,matrix}
\usepackage{bbm}
\usepackage{authblk}
\usepackage{subcaption}
\usepackage{xstring}
\usepackage{xparse}
\usepackage{comment}
\renewcommand{\Affilfont}{\footnotesize}

\usepackage{stmaryrd}
\usepackage{comment}
\usepackage{showlabels}
\usepackage{soul}

\usetikzlibrary{calc,patterns,shapes.geometric}
\usetikzlibrary{arrows,shapes,positioning}
\usetikzlibrary{decorations.markings,decorations.pathmorphing,decorations.pathreplacing}
\usetikzlibrary{decorations.text,decorations.shapes}

%%%%%%%% MACROS
\input{GempicLatexMacros}

\title{Coding Style and Conventions}

\author[1]{The Gempic development team}
\affil[1]{Max-Planck-Institut f\"ur Plasmaphysik, Garching, Germany}


\begin{document}

\maketitle

\begin{abstract}
This note describes some style conventions that should be followed in writing the GEMPIC code as well as the usage of the linting and formatting tools \verb|clang-tidy| and \verb|clang-format|.
\end{abstract}

\tableofcontents

\section{Conventions}
GEMPIC uses a set of style conventions expanded from \href{https://warpx.readthedocs.io/en/latest/developers/contributing.html#style-and-conventions}{the WarpX project}. The rationale behind some of the conventions are explained there.
These conventions only apply to source and header files. Checks that cannot be automatically enforced are marked with (*). Everything else could be fixed by running \verb|scripts/sanitize_with_python.py|.

\subsection{General conventions}
\begin{itemize}
    \item Four spaces for indentation (no tabs)
    \item <= 100 characters per line
    \item Spaces around assignment operator (\lstinline{=})
    \item Spaces after the function name for function definitions, \emph{e.g.} \lstinline{my_function ()}. Note that there should be no such spaces when simply calling or declarating a function.
    \item Overriding virtual functions are declared without the \lstinline{virtual} keyword but with \lstinline{override} or \lstinline{final}
    \item Function or class definitions have the starting \verb|{| brace on a new line
    \item \lstinline{#include}s are surrounded by quotation marks (\verb|"..."|) for GEMPIC headers and angular brackets (\verb|<...>|) otherwise
    \item (*) \emph{Do not} put \lstinline{using namespace} statements in header (\verb|.H|) files
    \item (*) \emph{Do not} put third party \lstinline{using namespace} statements in source (\verb|.cpp|) files
\end{itemize}
\subsection{Naming conventions}
\begin{itemize}
    \item Names for classes and namespaces are \verb|CamelCase|
    \item Names for functions, including member functions, are \verb|snake_case|
    \item Variable names are \verb|camelBack|. Exceptions are single character names optionally followed by \lstinline{x}, \lstinline{y} or \lstinline{z}, \emph{e.g.} \verb|E|, \lstinline{Bx} or \lstinline{Hz}.
    \item Template variable names are \verb|CamelCase| if they're type names and \verb|camelBack| if not
    \item Macros are \lstinline{UPPER_CASE}
    \item (Non-static) class member variables have the \lstinline{m_} prefix
    \item Static class variables have the \lstinline{s_} prefix
    \item Folder names are \verb|CamelCase|
    \item File names are prefixed with \verb|GEMPIC_| and then in \verb|CamelCase|, \emph{e.g.} \verb|GEMPIC_SomeModule.H|
    \item Test file names are \verb|CamelCase| and then postfixed with \verb|_test|, \emph{e.g.} \verb|SomeModule_test.cpp|
    \item (*) Namespaces follow the outer folder structure, such that code in \verb|Src/Io/Parameters| would be in the scope \lstinline{Gempic::Io}
    \item (*) Implementation details internal to a module can be "hidden" in an \verb|Impl| namespace, \emph{e.g.} \verb|Gempic::Io::Impl|
\end{itemize}
Exceptions to these conventions are unavoidable when inheriting from or calling third party libraries, but otherwise they should be followed and will be enforced with automatic tools.

\section{Automatic Linting and Formatting}
To help enforce the above conventions we use two clang extra tools: \verb|clang-format| and \verb|clang-tidy|.

\verb|clang-format| is mostly responsible for formatting whitespace: \verb|{| on newlines, spaces around assignment operators, order of \lstinline{#include} statements, indentation, etc. It can be run in seconds.

\verb|clang-tidy| is responsible for semantic style choices: parentheses around macro arguments, no cstyle casting, correct case for variables, functions and classes, etc. It takes minutes to run, with recent versions taking longer.

\begin{warning}
    Unfortunately, these tools still have some bugs. It is therefore recommended not to run them on their own, but instead run the script \verb|scripts/sanitize_with_python| which speeds up the runtime of \verb|clang-tidy| and attempts to correct mistakes incurred by both tools.
    To format and lint your code before committing run:\\
    \verb|python scripts/sanitize_with_python.py|\\
    The formatting should be done with clang-format and clang-tidy version 14-18.
\end{warning}
\begin{tip}
    It's highly recommended to run the sanitize script \emph{after} configuring compilation for 3D and \emph{after} staging your changes.
\end{tip}

\subsection*{Frequently Asked Questions: Automatic Linting and Formatting}
\begin{itemize}
    \item \emph{\ldots why?}

    Clean code is easier to read and debug. Strict style conventions provide a safety net that's hopefully not too restrictive.

    \item \emph{Your style conventions are too restrictive. Can I not just focus on programming?}
    
    Sure you can! Just remember to run \verb|scripts/sanitize_with_python.py| when you're done (and check that it still compiles and runs).

\end{itemize}

\end{document}
